% ===================================================================
%  TAFEL Nr. 5 — Das Haus und wie es gebaut wird.
%  Traduction 2026 d'apres texto/io, controlee sur texto/fr et
%  texto/en. Consigne : voir texto/de/10-tafel-01.tex.
%
%  Le tableau 5 est un DIALOGUE, et les deux volumes composent le nom
%  du parleur en petites capitales : on garde \textsc{}, que le rendu
%  laisse tomber en gardant le texte.
%
%  Les renvois du plan sont des LETTRES — (a) le couloir, (b) le
%  salon... — et non des numeros. On les ecrit comme l'ido les ecrit,
%  y compris la ou l'imprimeur a compose « (1) » pour « (l) ».
%
%  C'EST LA SEULE PAGE DES SEIZE OU L'IDO PARLE ALLEMAND EN TOUTES
%  LETTRES, ET LA SEULE OU LA COLONNE ALLEMANDE PUISSE DIRE QUE
%  GUIGNON S'EST TROMPE.
%
%  La note de l'alinea 25 propose une racine neuve, « balk-o », et
%  l'appuie sur l'allemand : « Balk-o = deutsch Balken, englisch
%  joist, französisch solive... » Elle precise ce qu'elle veut dire :
%  une solive, « qui n'est ni une trabo (francais poutre) ni une
%  petite poutre ».
%
%  OR « BALKEN » EST EXACTEMENT LA POUTRE QU'IL EXCLUT. L'allemand a
%  cinq mots pour cette charpente, et ils ne se confondent pas :
%     Balken      la poutre porteuse — le « trabo » qu'il refuse ;
%     Deckenbalken la solive de plancher — le mot qu'il cherchait ;
%     Sparren     le chevron du toit ;
%     Pfette      la panne, la piece horizontale du comble ;
%     Riegel      la traverse du pan de bois.
%  Guignon a pris LE MOT LE PLUS GENERAL DE LA LANGUE LA PLUS RICHE
%  POUR NOMMER LA CHOSE LA PLUS PARTICULIERE. La colonne allemande
%  ecrit « Deckenbalken » a l'alinea 25, parce que c'est de cela que
%  la planche parle ; la note, elle, cite « Balken », parce que c'est
%  le mot que la note cite. La source ne bouge pas.
%
%  ET LE RESTE DE LA PAGE DIT POURQUOI IL N'AVAIT PAS LE CHOIX : les
%  metiers du batiment sont le vocabulaire le plus purement allemand
%  des seize tableaux, et il est vieux de sept siecles.
%     Maurer, Zimmermann, Dachdecker, Schreiner, Glaser, Schlosser,
%     Schmied, Maler, Gipser, Tapezierer, Gärtner, Steinmetz,
%     Handlanger, Lehrling, Geselle, Meister.
%  Ce sont des noms de CORPORATION, fixes par les statuts des villes
%  au XIIIe siecle, et pas un n'a bouge. En face, les metiers d'apres
%  la corporation sont tous latins ou francais : Architekt,
%  Bauunternehmer, Ingenieur, Elektriker,
%  Plan. LA DATE DU METIER SE LIT DANS LA FORME DE SON NOM,
%  et l'electricien de l'alinea 41 est de 1890 comme le macon de
%  l'alinea 13 est de 1250.
%
%  UN MOT DONT LA COLONNE AUTRICHIENNE AURA BESOIN, ET QU'ON SIGNALE
%  ICI : le zingueur de l'alinea 28 est « Klempner » au nord,
%  « Spengler » au sud et en Autriche, « Flaschner » en Souabe ; le
%  menuisier de l'alinea 31 est « Schreiner » a l'ouest et
%  « Tischler » a l'est. Cette colonne-ci ecrit l'allemand standard
%  d'Allemagne — Klempner, Schreiner. Le partage n'est pas une faute a
%  corriger : c'est la matiere de la colonne autrichienne, quand elle
%  viendra.
%
%  DEUX INVERSIONS PREVUES ET EVITEES, la faute du tableau 4 : le
%  complement allemand remonte devant le verbe. Les alinéas 12 et 21
%  sont les plus charges de la page — huit renvois chacun — et les
%  deux ont demande qu'on rejette le complement derriere son groupe.
% ===================================================================

%%K t05-apar-1 apar
\VUsaut{6.0mm}
\VUtitre{15.0pt}{60}{TAFEL Nr. 5}
\VUsaut{1.4mm}
\VUtitre{11.4pt}{0}{Das Haus und wie es gebaut wird.}
\VUsaut{2.2mm}

%%K t05-tit-1 sub
\VUpk{10.2pt}{40}{Eine glückliche Begegnung.}
\VUsaut{3.0mm}

%%K t05-01-1 p
\textsc{Johann}. --- Onkel, ich wüsste zu gern, wie ein \VUgras{Haus} gebaut wird.

%%K t05-01-2 p
\textsc{Der Onkel} \textsuperscript{(80)}. --- Das ist nicht schwer, mein Junge. Komm
mit, ich zeige dir das Haus, das Herr Bernhard \textsuperscript{(79)} bauen lässt
und in das ich einziehen werde.

%%K t05-01-3 p
\textsc{Johann}. --- Und wer hat den \VUgras{Plan} \textsuperscript{(76)} für dieses
Haus gezeichnet?

%%K t05-01-4 p
\textsc{Der Onkel}. --- Der \VUgras{Architekt} \textsuperscript{(75)}; er hat ihn sehr
klar gezeichnet, er wurde genehmigt, und der \VUgras{Bauunternehmer}
\textsuperscript{(77)} hat mit der Arbeit angefangen.

%%K t05-01-5 p
\textsc{Johann}. --- Und wer ist der Unternehmer?

%%K t05-01-6 p
\textsc{Der Onkel}. --- Herr Weiß, der Vater deines Freundes Jakob. Er führt die
Arbeiter zusammen mit Herrn Roux, dem Architekten.

%%K t05-01-7 p
Sieh mal an! Da kommt Herr Weiß gerade im richtigen Augenblick mit seinem
\VUgras{Meterstab} \textsuperscript{(78)}, und Herr Roux mit seinem Plan.

%%K t05-01-8 p
Guten Tag, die Herren. Wie geht es Ihnen?

%%K t05-01-9 p
\textsc{Herr Roux}. --- Danke, sehr gut. Guten Tag, Johann; wie der Junge gewachsen
ist!

%%K t05-01-10 p
\textsc{Der Onkel}. --- Nicht wahr? Er ist schon ein kleiner Mann. Er ist sehr
ernst; alles, was er sieht, will er erklärt haben. Heute möchte er wissen, wie ein
Haus gebaut wird.

%%K t05-tit-2 sub
\VUpk{10.2pt}{40}{Die Arbeiter.}
\VUsaut{3.0mm}

%%K t05-01-11 p
\textsc{Herr Weiß}. --- Dann komm mit, ich erkläre es dir gleich. Ich mag Kinder,
die etwas lernen wollen.

%%K t05-01-12 p
Siehst du den Arbeiter mit der \VUgras{Schaufel} \textsuperscript{(82)} in der Hand:
das ist ein \VUgras{Erdarbeiter} \textsuperscript{(81)}. Er hebt einen
\VUgras{Graben} \textsuperscript{(46)} für die Wasserleitung aus. Er hat auch den
Boden ausgehoben, auf dem das Haus steht, damit der Maurer die vier
\VUgras{Mauern} hochziehen konnte. Der Teil dieser Mauern, der im Boden steckt,
heißt das \VUgras{Fundament} \textsuperscript{(2)}. Der Raum, den das Fundament
umschließt, bildet den \VUgras{Keller} \textsuperscript{(3)}. Dieser Keller hat ein
kleines Fenster, das \VUgras{Luftloch} \textsuperscript{(5)} heißt, eine
\VUgras{Tür} \textsuperscript{(4)} und eine Treppe.

%%K t05-01-13 p
Der \VUgras{Maurer} \textsuperscript{(108)} hat dann die Mauern weiter hochgezogen und
Öffnungen für die \VUgras{Türen} \textsuperscript{(8)} und die \VUgras{Fenster}
\textsuperscript{(17)} gelassen.

%%K t05-01-14 p
\textsc{Johann}. --- Aber ich sehe ihn ja gar nicht, diesen Maurer!

%%K t05-01-15 p
\textsc{Herr Weiß}. --- Am Haus ist er fertig. Er ist drüben beim \VUgras{Stall}
\textsuperscript{(54)}, oben auf seinem \VUgras{Gerüst} \textsuperscript{(110)}, und
mauert die Wand des \VUgras{Waschhauses} \textsuperscript{(56)}.

%%K t05-01-16 p
Er hat eine Kelle, einen Mörtelkasten und ein \VUgras{Senklot}
\textsuperscript{(109)}. Aber diese Namen wirst du nicht behalten.

%%K t05-01-17 p
\textsc{Johann}. --- Doch, bestimmt, denn ich werde meinen Onkel um eine kleine
Kelle und einen Mörtelkasten bitten, und ein Senklot mache ich mir selbst aus
einem Stück Schnur.

%%K t05-tit-3 sub
\VUsaut{3.0mm}
\VUpk{10.2pt}{40}{Die Baustoffe.}
\VUsaut{3.0mm}

%%K t05-01-18 p
\textsc{Der Onkel}. --- Sehr gut. Aber sag mal, Johann, was braucht man, um ein
Haus zu bauen?

%%K t05-01-19 p
\textsc{Johann}. --- Man braucht \VUgras{Hausteine} \textsuperscript{(59)} und
\VUgras{Mörtel} \textsuperscript{(65)}.

%%K t05-01-20 p
\textsc{Der Onkel}. --- Genau. Sieh dir den Mann mit dem großen Hut auf seinem
\VUgras{Karren} \textsuperscript{(136)} an. Das ist ein \VUgras{Fuhrmann}
\textsuperscript{(135)}. Jeden Tag bringt er \VUgras{Steinblöcke}
\textsuperscript{(60)} hierher. Er lädt seinen Karren im Hof ab, und die
\VUgras{Steinmetze} \textsuperscript{(83)} behauen die Blöcke und schneiden sie mit
ihrer \VUgras{Steinsäge} \textsuperscript{(85)}. Ein \VUgras{Handlanger}
\textsuperscript{(146)} trägt sie zum Maurer, der sie versetzt.

%%K t05-01-21 p
Unterdessen rührt der \VUgras{Mörtelmischer} \textsuperscript{(87)} den Mörtel an. Er
braucht \VUgras{Sand} \textsuperscript{(62)}, \VUgras{Kalk} \textsuperscript{(63)} und
\VUgras{Wasser} \textsuperscript{(64)}. Der Kalk steht neben ihm in einem
\VUgras{Sack} \textsuperscript{(71)}; ein \VUgras{Maurergeselle}
\textsuperscript{(111)} bringt ihm den Sand in einer \VUgras{Schubkarre}
\textsuperscript{(112)}, und der \VUgras{Lehrling} \textsuperscript{(113)} steht an
der Pumpe \textsuperscript{(58)}. Er füllt einen \VUgras{Eimer}
\textsuperscript{(114)}. Der Mörtel ist bald fertig. Der \VUgras{Mörtelträger}
\textsuperscript{(107)} nimmt ihn und trägt ihn die Leiter hinauf auf das Gerüst, wo
ein Maurer diese Seite des Hauses fertig macht.

%%K t05-01-22 p
Und jetzt, wie heißt der Arbeiter, der das \VUgras{Dach} \textsuperscript{(31)}
macht?

%%K t05-01-23 p
\textsc{Johann}. --- Das ist der \VUgras{Dachdecker} \textsuperscript{(118)}.

%%K t05-tit-4 sub
\VUsaut{3.0mm}
\VUpk{10.2pt}{40}{Das Dach des Hauses.}
\VUsaut{3.0mm}

%%K t05-01-24 p
\textsc{Herr Weiß}. --- Ja, aber zuerst kommt der \VUgras{Zimmermann}
\textsuperscript{(115)}.

%%K t05-01-25 p
Der Zimmermann macht den \VUgras{Dachstuhl} \textsuperscript{(35)}. Er verbindet die
\VUgras{Balken} \textsuperscript{(37)} mit \VUgras{Holznägeln}
\textsuperscript{(117)}. Die Arbeiter da oben, in ihren weiten Samthosen, sind
Zimmerleute. Der eine hält einen kleinen Balken, und der andere schlägt mit seinem
\VUgras{Beil} \textsuperscript{(116)} einen Holznagel zurecht. Der neben dem
\VUgras{Schornstein} \textsuperscript{(41)} ist der Dachdecker. Er nagelt Latten auf
die \VUgras{Deckenbalken} (*) und \VUgras{Schiefer} \textsuperscript{(66)} auf die
Latten. Manchmal deckt er auch mit Ziegeln.

%%K t05-noto-1 noto
\VUnotes{143.2mm}{%
(*) \VUgras{Balk-o} = deutsch \textit{Balken}, englisch \textit{joist},
französisch \textit{solive}, italienisch \textit{travicello}, russisch
\textit{balka}, spanisch und portugiesisch \textit{viga}. Eine Fachwurzel, die
noch nicht übernommen ist, hier aber vorgeschlagen wird, um den Sinn des
französischen Wortes \textit{solive} wiederzugeben, das weder ein \textit{trabo}
(französisch \textit{poutre}) noch ein kleiner Balken ist. Es ist ein
zugerichtetes Holz, das einen Fußboden und eine Decke trägt.%
}

%%K t05-01-26 p
Das Dach des Stalls ist mit \VUgras{Ziegeln} \textsuperscript{(55)} gedeckt, das des
Hauses mit \VUgras{Schiefer} \textsuperscript{(32)}, und das der Veranda ist aus
\VUgras{Zink} \textsuperscript{(24)}.

%%K t05-01-27 p
\textsc{Johann}. --- Macht der Dachdecker auch die Zinkdächer?

%%K t05-01-28 p
\textsc{Herr Weiß}. --- Nein, das ist der \VUgras{Zinkarbeiter}
\textsuperscript{(121)} oder der \VUgras{Klempner}, der da ein
\VUgras{Dachfenster} \textsuperscript{(38)} in das Dach des Hauses einsetzt. Er legt
auch die \VUgras{Dachrinnen} \textsuperscript{(43)} und die \VUgras{Fallrohre}
\textsuperscript{(42)}. Er lötet das Zink und das Weißblech. Er ist es auch, der den
\VUgras{Blitzableiter} \textsuperscript{(40)} und die \VUgras{Wetterfahne}
\textsuperscript{(39)} auf dem \VUgras{Giebel} \textsuperscript{(36)} befestigt hat.

%%K t05-01-29 p
\textsc{Johann}. --- Ist das Haus dann fertig?

%%K t05-tit-5 sub
\VUsaut{3.0mm}
\VUpk{10.2pt}{40}{Das Werkzeug.}
\VUsaut{3.0mm}

%%K t05-01-30 p
\textsc{Herr Weiß}. --- Nicht ganz. Der \VUgras{Gipser} \textsuperscript{(99)} mauert
mit \VUgras{Ziegeln} \textsuperscript{(61)} die Wände, die es in die einzelnen
Zimmer teilen. Er trägt \VUgras{Gips} \textsuperscript{(70)} auf die
\VUgras{Decken} \textsuperscript{(26)} und auf die Wände auf. Gerade jetzt arbeitet
er an der Decke der \VUgras{Veranda} \textsuperscript{(33)}. Wie der Maurer hat er
eine \VUgras{Kelle} \textsuperscript{(100)} und einen \VUgras{Mörtelkasten}
\textsuperscript{(101)}.

%%K t05-01-31 p
Dann macht der Schreiner die Türen, die Fenster, die \VUgras{Fußböden}
\textsuperscript{(16)} und die \VUgras{Fensterläden} \textsuperscript{(19)}.

%%K t05-01-32 p
Gerade jetzt arbeitet er im Hof an seiner \VUgras{Hobelbank}
\textsuperscript{(90)}. Er hat eine \VUgras{Säge} \textsuperscript{(94)}, um das
\VUgras{Holz} \textsuperscript{(69)} zu schneiden, einen \VUgras{Hobel}
\textsuperscript{(91)}, eine \VUgras{Schraubzwinge} \textsuperscript{(92)}, ein
\VUgras{Stemmeisen} \textsuperscript{(94 \textit{bis})}, einen \VUgras{Zirkel}
\textsuperscript{(95)} und einen hölzernen \VUgras{Klüpfel}
\textsuperscript{(95 \textit{bis})}.

%%K t05-01-33 p
\textsc{Johann}. --- Oh, Schreiner ist ja noch schöner als Maurer. Onkel, kaufst du
mir eine Hobelbank, eine Säge und einen Hobel? Ich will eine Hundehütte für Medor
bauen.

%%K t05-01-34 p
\textsc{Der Onkel}. --- Ja, mein Junge.

%%K t05-01-35 p
\textsc{Johann}. --- Danke schön! Und wer ist der Mann da vorn an der Haustür?

%%K t05-tit-6 sub
\VUsaut{3.0mm}
\VUpk{10.2pt}{40}{Der Anstrich.}
\VUsaut{3.0mm}

%%K t05-01-36 p
\textsc{Herr Weiß}. --- Das ist ein \VUgras{Maler} \textsuperscript{(102)}. Er steht
auf seiner Leiter mit einem Topf \VUgras{Farbe} \textsuperscript{(103)} und seinem
\VUgras{Pinsel} \textsuperscript{(104)}. Er streicht das Holz der Haustür an, damit
es länger hält.

%%K t05-01-37 p
Er ist es auch, der in den Zimmern die Tapeten klebt.

%%K t05-01-38 p
Der \VUgras{Tapezierer} \textsuperscript{(107)} auf einer \VUgras{Leiter}
\textsuperscript{(106)}, draußen auf dem \VUgras{Balkon} \textsuperscript{(28)},
bringt eine \VUgras{Markise} \textsuperscript{(29)} an.

%%K t05-01-39 p
Der Arbeiter am großen Fenster ist der \VUgras{Glaser} \textsuperscript{(96)}. Er
setzt die \VUgras{Scheiben} \textsuperscript{(18)} mit \VUgras{Kitt}
\textsuperscript{(73)} ein. Der andere Arbeiter, der an der Kellertür kniet, ist ein
\VUgras{Schlosser} \textsuperscript{(97)}. Er setzt ein \VUgras{Schloss}
\textsuperscript{(6)} ein. Seine \VUgras{Feile} \textsuperscript{(98)} liegt neben
ihm.

%%K t05-01-40 p
\textsc{Johann}. --- Ist der Mann, der da mit seinem \VUgras{Hammer}
\textsuperscript{(84)} auf ein Stück Eisen schlägt, auch ein Schlosser?

%%K t05-01-41 p
\textsc{Herr Weiß}. --- Nicht ganz --- er ist ein \VUgras{Schmied}
\textsuperscript{(125)}. Er schmiedet \VUgras{Eisenwerk} \textsuperscript{(67)} für
den \VUgras{Elektriker} \textsuperscript{(124)}, der auf dem Dach der
\VUgras{Remise} \textsuperscript{(51)} steht. Er hat eine \VUgras{Esse}
\textsuperscript{(126)}, einen \VUgras{Amboss} \textsuperscript{(127)} und eine
\VUgras{Zange} \textsuperscript{(128)}.

%%K t05-01-42 p
Der Elektriker befestigt den \VUgras{Kupferdraht} \textsuperscript{(44)} des
Telefons.

%%K t05-tit-7 sub
\VUpk{10.2pt}{40}{Die Gartenarbeit.}
\VUsaut{3.0mm}

%%K t05-01-43 p
\textsc{Der Onkel}. --- Hinten im \VUgras{Hof} \textsuperscript{(45)}, neben dem
\VUgras{Gitter} \textsuperscript{(47)} und dem \VUgras{Einfahrtstor}
\textsuperscript{(48)}, siehst du den \VUgras{Gärtner} \textsuperscript{(129)}. Er
macht den kleinen \VUgras{Garten} \textsuperscript{(49)} fertig. Er hat die Erde mit
seinem \VUgras{Spaten} \textsuperscript{(131)} umgegraben, er sät und harkt mit dem
\VUgras{Rechen} \textsuperscript{(130)}. Dann wird er mit seiner
\VUgras{Gießkanne} \textsuperscript{(132)} die \VUgras{Beete}
\textsuperscript{(150)} gießen, und die Pflanzen werden wachsen.

%%K t05-01-44 p
Der \VUgras{Stallknecht} \textsuperscript{(133)}, der gerade das Pferd versorgt und
gefüttert hat, putzt den \VUgras{Wagen} \textsuperscript{(52)} mit einem Schwamm,
einer \VUgras{Handpumpe} \textsuperscript{(134)} und einem Eimer Wasser. Dann wird er
ihn in die Remise stellen und die Tür mit dem schweren \VUgras{Riegel}
\textsuperscript{(53)} verschließen.

%%K t05-01-45 p
So, jetzt bist du hoffentlich zufrieden; Erklärungen hast du genug bekommen.

%%K t05-01-46 p
\textsc{Johann}. --- Ja, Onkel, aber ich würde das Haus gern auch von innen sehen.

%%K t05-01-47 p
\textsc{Der Onkel}. --- Das hat Zeit bis morgen. Herr Roux wird dir den Plan
erklären und dir sagen, wie jedes Zimmer heißt.

%%K t05-tit-8 sub
\VUsaut{3.0mm}
\VUpk{10.2pt}{40}{Die Einteilung des Hauses.}
\VUsaut{2.4mm}

%%K t05-apar-2 apar
{\centering\textit{(Siehe den Plan.)}\par}
\VUsaut{2.4mm}

%%K t05-01-48 p
\textsc{Der Architekt}. --- Komm, Johann, ich führe dich durch die verschiedenen
Teile des Hauses.

%%K t05-01-49 p
Gehen wir ins \VUgras{Erdgeschoss} \textsuperscript{(7)}. Hier ist der Knopf der
\VUgras{Klingel} \textsuperscript{(11)}. Wir steigen die drei \VUgras{Stufen}
\textsuperscript{(14)} der \VUgras{Freitreppe} \textsuperscript{(9)} hinauf, wir
überschreiten die \VUgras{Schwelle} \textsuperscript{(10)} der \VUgras{Haustür}
\textsuperscript{(8)}, und schon stehen wir am Fuß der \VUgras{Treppe}
\textsuperscript{(13)}.

%%K t05-01-50 p
\textsc{Johann}. --- Das ist der Gang.

%%K t05-01-51 p
\textsc{Der Architekt}. --- Nein, das ist die \VUgras{Diele} \textsuperscript{(12)}.
Der \VUgras{Gang} \textsuperscript{(a)} ist ganz hinten. Rechts liegen der
\VUgras{Salon} \textsuperscript{(b)} und das \VUgras{Esszimmer}
\textsuperscript{(c)}; dem Esszimmer gegenüber die \VUgras{Küche}
\textsuperscript{(d)} und die \VUgras{Anrichte} \textsuperscript{(e)}; dann, nach
vorn hinaus, das \VUgras{Arbeitszimmer} \textsuperscript{(f)}. Die
\VUgras{Veranda} \textsuperscript{(23)} mit ihrer \VUgras{Glastür}
\textsuperscript{(25)} liegt neben dem Salon.

%%K t05-01-52 p
Jetzt gehen wir die Treppe hinauf. Halt dich am \VUgras{Geländer}
\textsuperscript{(15)} fest und zähl die Stufen. Es sind vierzehn. Hier ist der
\VUgras{Treppenabsatz} \textsuperscript{(m)}: wir sind im ersten \VUgras{Stock}
\textsuperscript{(27)}.

%%K t05-tit-9 sub
\VUsaut{3.0mm}
\VUpk{10.2pt}{40}{Der erste Stock.}
\VUsaut{3.0mm}

%%K t05-01-53 p
Der Treppe gegenüber liegen der \VUgras{Abort} \textsuperscript{(20)} und das
\VUgras{Ankleidezimmer} \textsuperscript{(j)}. Rechts ein schönes
\VUgras{Schlafzimmer} \textsuperscript{(g)} mit zwei Fenstern zur
\VUgras{Vorderseite} \textsuperscript{(1)} des Hauses. Links das
\VUgras{Badezimmer} \textsuperscript{(1)} und das \VUgras{Gästezimmer}
\textsuperscript{(i)} mit einem großen \VUgras{Alkoven} \textsuperscript{(h)}. Neben
dem Schlafzimmer liegt das \VUgras{Kinderzimmer} \textsuperscript{(k)}. Es geht auf
eine breite \VUgras{Terrasse} \textsuperscript{(21)} hinaus. Unter dieser Terrasse
liegt ein zweiter Abort \textsuperscript{(20)}, und an der Wand hängt die
\VUgras{Glocke} \textsuperscript{(22)}, die zum Essen läutet.

%%K t05-01-54 p
\textsc{Johann}. --- Gehen wir hinauf in den \VUgras{zweiten Stock}.

%%K t05-01-55 p
\textsc{Der Architekt}. --- Einen zweiten Stock gibt es nicht. Unter dem Dach sind
nur \VUgras{Dachkammern} \textsuperscript{(33)} und der Speicher
\textsuperscript{(34)}. Das ist der ganze Rundgang --- wir können wieder
hinuntergehen.

%%K t05-01-56 p
\textsc{Johann}. --- Danke, mein Herr, das war wirklich interessant. Ich erzähle
meinem Vater alles, was ich gesehen habe.
\VUsaut{4.0mm}
\VUfilet{22mm}
